\documentclass[a4paper,11pt,titlepage]{article}
\usepackage{amssymb}
\usepackage{hyperref}
\usepackage{booktabs}
\usepackage{geometry}
\usepackage{listings}
\usepackage{graphicx}
\usepackage{natbib}
\usepackage{todonotes}
\usepackage{enumitem}
\usepackage{multirow}
\usepackage{tikz}
\usepackage{fancyhdr}
\usepackage{xcolor,colortbl}
\usepackage{eurosym}
\usetikzlibrary{arrows,decorations.pathmorphing,backgrounds,positioning,fit,petri,matrix,folding}

\setlength\headheight{20pt}
\addtolength\topmargin{-10pt}
\addtolength\footskip{20pt}

\fancypagestyle{plain}{%
\fancyhf{}
\fancyfoot[RO,LE]{\sffamily\bfseries\thepage}
\renewcommand{\headrulewidth}{0pt}
\renewcommand{\footrulewidth}{1pt}
}


\newcommand{\HRule}{\rule{\linewidth}{0.5mm}}

\newcommand{\uni}{Eindhoven University of Technology}
\newcommand{\vak}{Database Technology}
\newcommand{\vakcode}{2ID35}
\newcommand{\essaytitle}{Smooth Scan Analysis}
\newcommand{\stad}{Eindhoven}
\newcommand{\tbsep}{\ \ \textbar \ \textbar \ \textbar \ \textbar \ \textbar \ \textbar \ \textbar \ \textbar \ \textbar \ \textbar \ \ }

\pagestyle{fancy}
\fancyhf{}
\fancyhead[RO,LE]{\vak}
\fancyhead[LO,RE]{\uni}
\fancyfoot[RO,LE]{\sffamily\bfseries\thepage}
\fancyfoot[LO,RE]{\essaytitle}
\renewcommand{\headrulewidth}{1pt}
\renewcommand{\footrulewidth}{1pt}

\bibliographystyle{plain}
\hypersetup{pdfborder={0 0 0}}

\setlength{\parskip}{10pt}

\geometry{
	includeheadfoot,
	margin=2.54cm
}
\author{
	Nicky Advokaat (0740567) - \texttt{n.advokaat@student.tue.nl}
	\and
	Robbert Jongeling (0747896) - \texttt{r.m.jongeling@student.tue.nl}
	\and
	Bram Kohl (0746107) - \texttt{b.j.e.kohl@student.tue.nl}
	\and
	Jasper Selman (0741516) - \texttt{j.w.m.selman@student.tue.nl}
	\and
	Ramon de Vaan (0758873) - \texttt{r.d.vaan@student.tue.nl}
}
\date{\today}

\begin{document}
	\begin{center}

% Upper part of the page
		\includegraphics[width=0.15\textwidth]{images/tuelogo}\\[1cm]

		\textsc{\LARGE \uni}\\[1.6cm]


        \textsc{\LARGE \vak}\\[0.5cm]

% Title
\HRule \\[0.4cm]
{ \huge \bfseries \essaytitle}\\[0.4cm]

\HRule \\[1.5cm]

% Author and supervisor
	\emph{Group TODO:}\\
    \begin{tabular}{l l l}
	Sander \textsc{Breukink} & 0741209 & \href{mailto:s.c.breukink@student.tue.nl}{\texttt{s.c.breukink@student.tue.nl}}\\
	Rianne \textsc{Conijn} & 0740635 & \href{mailto:m.a.conijn@student.tue.nl}{\texttt{m.a.conijn@student.tue.nl}}\\
	Harm \textsc{van Schaaijk} & TODO & \href{mailto:h.a.h.v.schaaijk@student.tue.nl}{\texttt{h.a.h.v.schaaijk@student.tue.nl}}\\
	Jasper \textsc{Selman} & 0741516 & \href{mailto:j.w.m.selman@student.tue.nl}{\texttt{j.w.m.selman@student.tue.nl}}\\
	Tom \textsc{Vogels} & 0871231 & \href{mailto:t.j.h.vogels@student.tue.nl}{\texttt{t.j.h.vogels@student.tue.nl}}
    \end{tabular}
		\vfill

% Bottom of the page
{\large \today} \\
\stad

	\end{center}

    \newpage
\section{Context and motivation}
Nowadays big data is booming business. A way to handle all this data in a fast efficient manner is very important for this cause. That is why at this time there are a lot of query optimizers, but these optimizers need statistics about the big data to create good query plans. The problem is in many of these cases such statistics are sparse or even non-existent. \\
\\
This is why Borovica, Idreas, Ailamakki, Zukowsku, and Fraser (2015) \cite{smoothscan} deisgned a new method for the query optimizations decisions, the process of deciding which physical operators are used and in what order. Note that these decisions can affect the response time by several orders of magnitude and are therefore crucial in creating an optimally performing database. The proposed method is called \emph{Smooth Scan} and its results are compared to those of a \emph{full table scan}, an \emph{index scan} and a \emph{switch scan}. The \emph{smooth scan} is inspired by adaptive query processing, a technique which uses runtime feedback to modify query processing to get a better response time or more efficient CPU utilization (Deshpande, Ives, $\&$ Raman, 2007) \cite{query}.\\
\\ 
The motivation for this article is the severe impact several scanning methods have on the total execution time, a response time which, in this day and age, becomes increasingly important. It also seems that the cost models for performance estimations deviate increasingly as the system becomes more complex. Combined with the fact that the complexity of modern workloads and the technological shift towards cloud environments, query optimization becomes increasingly important. Because of these reasons there is enough motivation to discuss a new scanning method for query handling.

%The report also discusses the trade-off between CPU operations and I/O operations. Because it is possible to do several thousands of CPU operations in the time needed for one I/O operation, this trade-off is very often beneficial for the performance. Since the durations of these operations are subject to continuous change because of the fact hardware is always being improved, this too is a motivation for writing this article. 

\section{Research problem}
The main problem addressed in the paper is that once a decision is made by the query optimizer, this decision is fixed throughout the execution of a query. This can result in suboptimal plans and even a small estimation error can lead to drastically different performance results. These unpredictable results makes the system non-robust. The authors define robustness as: “the ability of a system to efficiently cope with unexpected and adverse conditions, and deliver near-optimal performance for all query inputs.” In their paper they try to achieve robust query processing.

\section{Smooth Scan}
In this section we briefly describe how the \emph{Smooth Scan} algorithm works. The former mentioned algorithms pick their strategy before the query is executed and stick to that strategy during the whole execution phase. The \emph{Smooth Scan} algorithm also picks a strategy at the start, but during the execution it might change strategy if necessary. \\
During the lifetime of \emph{Smooth Scan} the operator can be in three different modes. These modes are called \emph{Mode 1: Index Scan}, \emph{Mode 2: Entire Page Probe}, \emph{Mode 3: Flattening Access} and \emph{Mode 3+: Gradually Flattening Access}. In each of these modes the operator performs a gradually increasing amount of work as a result of the selectivity's increase. \\
In the first mode (in which \emph{Smooth Scan} always starts) a normal index scan is executed. Once the results cardinality threshold is exceeded it switches to mode 2. In this mode the operator performs a full table scan. If the result cardinality increases even more, \emph{Smooth Scan} changes to mode 3. In this mode the algorithm amortizes I/O cost over CPU cost (since you can do thousands of CPU operation during one I/O operation). This means the random selecting function is replaced with a sequential one. Mode 3+ is an expanding version of mode 3. The number of extra pages fetched for each single page it nees to access becomes larger and larger. Eventually this mode may change in a full table scan. For a more ellaborate explanation of the algorithm we refer to section 3 of the work of Borovica et al. \cite{smoothscan}.

\section{Claimed results}
The authors implemented Smooth Scan in PostgreSQL and claim that it achieves robust performance in a range of synthetic and real workloads, while being statistics-oblivious at the same time. Existing approaches fail to do so. This robust behaviour results in significant gains compared to when the original system makes a wrong decision, and only marginal overheads compared to when a correct decision can be made. Besides that, the authors also did the same tests with a SSD instead of a HDD. They claimed that SSD was even more benificial for \emph{Smooth Scan} than HDD.
%We can choose one (or more) of the statements below to check, as these provide some more explicit claims .
%- With Smooth Scan, PostgreSQL demonstrates robust performance. It does not suffer from extreme degradation and achieves good performance for all queries.
%- Smooth scan improves by a factor of 10 when compared to index scan  for selectivity 100$\%$.
%- Smooth scan on SSD is faster than full scan for selectivity lower than 20$\%$, and it is only 10$\%$ slower compared to full scan for 100$\%$ selectivity.

\section{Plan for verification}
In the current project we will try to verify some of the claims stated in the paper. 
Due to time constraints we will not implement the \emph{Smooth Scan}, but only the \emph{index scan}, the \emph{full table scan}, and the \emph{switch scan}.
For the \emph{index scan} a B\textsuperscript{+}-trees will be used. The \emph{switch scan}  will switch on a cardinality estimate, which still has to be set. This estimate will be comparable to the value set by Borovica et al. \cite{smoothscan}, who used a cardinality of 15K tuples out of 400M.
\\
These three scanswill be implented in PostgreSQL 9.2.1 DBMS and compared on execution time, for different selectivity levels. 
Hereby we will try to replicate Figure 12 from the paper \cite{smoothscan}.
All scans will be tested on a randomly generated database, in which multiple values occur more than once. 

%- ik vind het lastig om aan te geven welke claims we nu eigenlijk echt gaan verifiëren. Index scan met full table scan vergelijken is namelijk al best veel gedaan, dus we gaan dit natuurlijk gewoon verifiëren...

\section{Progress}

\begin{thebibliography}{9}

\bibitem{smoothscan}
	 Borovica, R., Idreos, S., Ailamaki, A., Zukowski, M., $\&$ Fraser, C.
	2015.	
 	Smooth Scan: One access path to rule them all.
	\emph{IEEE International conference on Data Engineering.}

\bibitem{query}
	Deshpande, A., Ives, Z., $\&$ Raman, V.
	2007.
	Adaptive query processing
	\emph{Foundations and Trends in Databases,}
	1(1), 1-140.
\end{thebibliography}

\end{document}
